\subsection{A quaternion surface quotient detecting the complete CM exponent}

\begin{result}{A.4.12}{Proved}{An actual characteristic surface detector of order $32\cdot3^{45}$}
Let $\Sigma=F(a,b)=\pi_1(E\setminus\{0\})$, with its displayed
oriented generators, and let
\[
 q_8:\Sigma\twoheadrightarrow Q_8,
 \qquad a\longmapsto i,\quad b\longmapsto j,
 \qquad N_8=\ker q_8.
                                                        \tag{QS.1}
\]
Put
\[
 L_8=\gamma_3(N_8)N_8^3,\qquad
 R_8=L_8\cap\ker\bigl(\Sigma\to(\mathbb Z/12)^2\bigr),
 \qquad \mathcal P_8=\Sigma/R_8.
                                                        \tag{QS.2}
\]
Both kernels are characteristic, including for continuous
automorphisms of the full profinite surface group.  The exact
orders are
\[
 |\Sigma/L_8|=8\cdot3^{45},\qquad
 |\mathcal P_8|=32\cdot3^{45}.
                                                        \tag{QS.3}
\]
The actual modular action on this surface quotient has an
explicit detected image $J$, with a specified homomorphism
\[
 \operatorname{im}\bigl(\operatorname{SL}_2(\mathbb Z)
             \to\operatorname{Out}(\mathcal P_8)\bigr)
       \twoheadrightarrow J
       \twoheadrightarrow\operatorname{SL}_2(\mathbb Z/12),
 \qquad |J|=1152\cdot3^8.
                                                        \tag{QS.4}
\]
Write $c$ for the image of the actual word $UL_0$, where
\[
 U=\begin{pmatrix}1&1\\0&1\end{pmatrix},\qquad
 L_0=\begin{pmatrix}1&0\\-1&1\end{pmatrix}.
\]
The centralizer $C_J(c)$ has order $6\cdot3^4$, and its exact
image modulo twelve is
\[
 \langle UL_0\bmod12\rangle\cong C_6.                    \tag{QS.5}
\]
Consequently the same image assertion holds for the centralizer
of $UL_0$ in the actual modular outer image on
$\mathcal P_8$: the detected quotient excludes the other thirty
commuting residues, and the six actual powers of $UL_0$ realize
all six remaining residues.  No injectivity of the first arrow
in (QS.4) is asserted or needed.

There is a literal common cover with the three-division surface.
For $K_3=\ker(\Sigma\to(\mathbb Z/3)^2)=\Pi_{1,9}$ put
$T_{73}=N_8\cap K_3$.  Then $T_{73}$ is free of rank $73$,
$[K_3:T_{73}]=8$, and $R_8\subseteq T_{73}$.  Thus
$\mathcal P_8$ contains the actual index-nine quotient
$K_3/R_8$ and its index-eight subgroup $T_{73}/R_8$.
This supplies a common finite surface target; it does not yet
supply the finite geometric fusion of all competing
three-division deck subgroups required in A.4.
\end{result}

\begin{proof}
\textbf{The characteristic quaternion map and its abelian join.}
Every generating pair of $Q_8$ consists of two anticommuting
noncentral elements of order four.  There are $6\cdot4=24$
such pairs, and a prescribed pair uniquely determines an
automorphism of $Q_8$: its product determines the image of $k$,
and the common square determines $-1$.  Hence all surjections
$F(a,b)\to Q_8$ have the same kernel.  Precomposition by any
automorphism permutes these surjections.  The identical argument
applies to continuous surjections from $\widehat\Sigma$.
This proves the claimed characteristic property of $N_8$;
the two operations in (QS.2) preserve it.

The index-eight free subgroup $N_8$ has rank nine.  Its free
class-two exponent-three quotient
\[
 V=N_8/L_8
   =\mathbb F_3^9\oplus\bigwedge^2\mathbb F_3^9
                                                        \tag{QS.6}
\]
has multiplication $(v,z)(w,t)=(v+w,z+t+2v\wedge w)$.
This is an explicit group law, with $2=1/2$ in $\mathbb F_3$;
it proves $|V|=3^{9+36}$ without using a class-three
Baker--Campbell--Hausdorff denominator.

The exponent-sum lattice of $N_8$ in $\Sigma^{\rm ab}=\mathbb Z^2$
is precisely $2\mathbb Z^2$.  Containment follows from
$Q_8^{\rm ab}=C_2^2$, and equality follows from the words
$a^2[a,b]^{-1}$ and $b^2[a,b]^{-1}$ in $N_8$.
The exponent lattice of $L_8$ is therefore $6\mathbb Z^2$.
Intersecting it with the kernel of abelian reduction modulo
twelve gives image $12\mathbb Z^2$, of index four in
$6\mathbb Z^2$.  The first isomorphism theorem applied to
$L_8\to(\mathbb Z/12)^2$ proves (QS.3).

Let $P$ denote the normal Sylow-three subgroup of
$\mathcal P_8$.  Its literal preimage in $\Sigma$ is
$N_8\cap\ker(\Sigma\to(\mathbb Z/4)^2)$.
Projection to $N_8/L_8$ identifies $P$ with $V$: the kernel
is $R_8$, and surjectivity follows because the image of $L_8$
modulo four is all of $2\mathbb Z^2/4\mathbb Z^2$.
In particular $P$ is characteristic.  The complementary
quotient, with no choice of splitting, is the actual group
\[
 \mathcal P_8/P
   =Q_8\mathbin{\times}_{C_2^2}(\mathbb Z/4)^2,
 \qquad |\mathcal P_8/P|=32.                              \tag{QS.7}
\]
The kernel $L_8/R_8$ of order four is central in
$\mathcal P_8$: for $g\in\Sigma$ and $h\in L_8$,
the commutator $[g,h]$ belongs to $L_8$ and has zero
exponent sums, hence belongs to $R_8$.  The image
$N_8/R_8$ is the product of $P$ with this central kernel,
by the established surjection $P\to N_8/L_8$.
Its conjugation on $P$ is consequently inner, so the
outer conjugation action factors through the actual
quotient $\mathcal P_8/(N_8/R_8)=Q_8$.
Conjugation of $\mathcal P_8$ on $P=V$, followed by passage to
outer automorphisms, has image the deck subgroup
$D\cong Q_8$.  Faithfulness on $V^{\rm ab}$ is checked by the
eight explicit matrices below.  Therefore restriction and the
abelianization give actual homomorphisms
\[
 \operatorname{Out}(\mathcal P_8)
 \longrightarrow
 N_{\operatorname{Out}(V)}(D)/D
       \ \times\ \operatorname{GL}_2(\mathbb Z/12).
                                                        \tag{QS.8}
\]
Indeed $\mathcal P_8^{\rm ab}=(\mathbb Z/12)^2$ by the same
exponent-lattice computation.  Changing an automorphism of
$\mathcal P_8$ by an inner automorphism changes its restriction
by an element of $D$ and acts trivially on that abelianization.
Define $J$ to be the image of the actual modular action under
(QS.8).  This proves that the first arrow in (QS.4) is an
actual quotient map before any enumeration is performed.

\textbf{The finite coordinates and the complete image calculation.}
The following Schreier basis fixes all markings.  Breadth-first
search in the $Q_8$ Cayley graph uses the letters
$a,b,a^{-1},b^{-1}$ in this order; positive edges are listed in
the order $1,i,j,k,-1,-i,-j,-k$, then $a,b$.  The nine
non-tree edges give
\[
\begin{split}
 (&bab a^{-1},\ b^2a^{-2},\ aba b^{-1},\ ab^2a,\ a^4,\\
  &a^2b^2,\ a^{-1}b^2a^{-1},\ b^{-1}ab^{-1}a^{-1},\
    ab^{-1}ab).
\end{split}                                               \tag{QS.9}
\]
Each word rewrites to its corresponding single basis letter;
this is checked in the retained producer.  On $\Sigma$ use
\[
 u:(a,b)\longmapsto(a,ba),\qquad
 l:(a,b)\longmapsto(ab^{-1},b),\qquad
 d_a=\operatorname{Ad}(a),\quad d_b=\operatorname{Ad}(b).
                                                        \tag{QS.10}
\]
Their closed-surface matrices are $U,L_0,I,I$.  Restrict these
four literal substitutions to (QS.9) and rewrite by the
Schreier algorithm.  These are the four explicit automorphisms
used in the certificate; their complete free words are saved.

For $W_0=\mathbb F_3^9$ the outer automorphisms of $V$ have the
coordinate description
\[
 \operatorname{Out}(V)
  =\bigl(\operatorname{Hom}(W_0,\bigwedge^2W_0)
          /\operatorname{ad}(W_0)\bigr)
       \rtimes\operatorname{GL}(W_0).
                                                        \tag{QS.11}
\]
The first vector space has dimension $324-9=315$.
In the ordered wedge basis $(i,j)$, $i<j$, the producer
collects a free word letter by letter using (QS.6).
An automorphism whose basis images have linear columns $M$
and central columns $B$ is represented by $(BM^{-1},M)$.
Row reduction of the nine maps $w\mapsto v\wedge w$ fixes
a specific complement for the inner maps.  Denote its projection
by $\pi$.  Multiplication and the action used throughout are
\[
 (C,M)(E,N)=(C+M\cdot E,MN),\qquad
 M\cdot E=\pi\bigl((\bigwedge^2M)EM^{-1}\bigr).
                                                        \tag{QS.12}
\]
All inverses and all row reductions are exact modulo three.

In these coordinates $d_a,d_b$ generate exactly eight
elements, distinct already by their linear matrices.  The
four generators normalize this deck subgroup, with equality
of both the linear and central coordinates.  The braid
relation for $u,l$ and the deck correction to $(ul)^6$ are
also verified in both coordinates.  Joining the four linear
matrices to their signed matrices modulo twelve gives a base
group of exactly $9216=8\cdot1152$ elements.  All of them are
enumerated by positive words in the four generators, retaining
the word for each element.  Quotienting by the eight deck
elements identifies this base with the full, explicitly
enumerated $\operatorname{SL}_2(\mathbb Z/12)$.

For each of the $9216$ base representatives $g_q$ and each of
the four generators $e$, form the actual Schreier element
$g_q e g_{qe}^{-1}$.  Its $315$-coordinate vector is retained.
These $36864$ vectors have rank eight.  With zero-based row
indices, an independent generating subset is
\[
 1356,\ 1384,\ 2000,\ 2028,\ 2141,\ 2153,\ 2293,\ 2337.
                                                        \tag{QS.13}
\]
Their span $W$ is invariant under each of the four actual
generators; the certificate checks this by another exact row
reduction.  Schreier generation now proves that $W$ is the
entire kernel of the joint base map.  The assertion is not
inferred from a sample of congruence generators: every edge
of the finite Cayley graph occurs.  Since the deck subgroup
has faithful linear image, $W$ meets it trivially.  Thus
\[
 1\longrightarrow W\cong C_3^8\longrightarrow J
 \longrightarrow\operatorname{SL}_2(\mathbb Z/12)
 \longrightarrow1,                                    \tag{QS.14}
\]
which proves the order in (QS.4).

\textbf{Every centralizer equation, including its deck correction.}
The matrix $R=UL_0\bmod12$ has $36$ commuting residues in
$\operatorname{SL}_2(\mathbb Z/12)$.  For each such residue $q$
choose the breadth-first positive $u,l$ word $g_q$ retained in
the certificate, and form its actual commutator with $c=ul$.
Write its coordinates as $(C_q,M_q)$.  Its linear component
is the component of a unique deck element $d_q=(E_q,M_q)$.
Every element of $J$ above $q$ is represented uniquely by
$(w,I)g_q$, $w\in W$, after quotienting by the deck subgroup.
By (QS.12) its commutator with $c$ has coordinates
\[
 \bigl(w+C_q-(M_qM_c)\cdot w,\ M_q\bigr).
\]
Consequently the exact centralizer condition in $J$ is
\[
 (I-M_qM_c)\cdot w=E_q-C_q,
       \qquad w\in W.                                  \tag{QS.15}
\]
This equation includes the actual deck element; setting it
silently to the identity is not part of the computation.
Both sides lie in the explicit eight-dimensional kernel
$W$.  Restriction to the eight pivot columns of its reduced
basis gives an $8\times8$ affine equation.  The certificate
retains all $36$ such matrices and right-hand sides, their
actual words $g_q,d_q$, and their four-coordinate cokernel
obstructions.  Each coefficient matrix has rank four.

For a shorter reproducible table, the quotient of the
$36$ commuting matrices by $\langle R\rangle$ has the following
six representatives.  They are uniquely specified by
$q_j\equiv I\pmod4$ and $q_j\equiv R^j\pmod3$.
The last column is the four-coordinate obstruction from the
retained row reduction of (QS.15):
\[
\begin{array}{c|c|c}
 j&q_j&\text{obstruction}\\\hline
 0&\left(\begin{smallmatrix}1&0\\0&1\end{smallmatrix}\right)
       &(0,0,0,0)\\
 1&\left(\begin{smallmatrix}9&4\\8&1\end{smallmatrix}\right)
       &(1,1,2,2)\\
 2&\left(\begin{smallmatrix}5&4\\8&9\end{smallmatrix}\right)
       &(0,1,2,2)\\
 3&\left(\begin{smallmatrix}5&0\\0&5\end{smallmatrix}\right)
       &(1,2,1,1)\\
 4&\left(\begin{smallmatrix}9&8\\4&5\end{smallmatrix}\right)
       &(0,2,1,1)\\
 5&\left(\begin{smallmatrix}1&8\\4&9\end{smallmatrix}\right)
       &(1,0,0,0)
\end{array}                                               \tag{QS.16}
\]
Multiplication by an actual power of $c$ bijects the lifting
sets in the same coset, so this table also proves the full
exclusion.  The retained calculation separately evaluates
all $36$ residues.  The six survivors are exactly the powers
of $R$, each with $3^4$ lifts.  This proves (QS.5) and the
claimed centralizer order.

\textbf{The literal common three-division cover.}
The maps $\Sigma\to Q_8$ and $\Sigma\to C_3^2$ give a
surjection to their product.  Indeed its image is a subdirect
product, and any common quotient of a group of order eight
and a group of order nine is trivial.  Its kernel is
$T_{73}=N_8\cap K_3$, of index $72$, hence of free rank $73$.
The exponent lattice $6\mathbb Z^2$ of $L_8$ shows directly
that $L_8\subseteq K_3$; therefore $R_8\subseteq T_{73}$.
All maps in the following square are induced by inclusions
of these displayed free subgroups, followed by the same
quotient $\Sigma\to\mathcal P_8$:
\[
\begin{array}{ccc}
 T_{73}/R_8&\lhook\joinrel\longrightarrow&N_8/R_8\\
 \big\downarrow&&\big\downarrow\\
 K_3/R_8&\lhook\joinrel\longrightarrow&\mathcal P_8.
\end{array}                                               \tag{QS.17}
\]
Their orders are respectively
$4\cdot3^{43},4\cdot3^{45},32\cdot3^{43},32\cdot3^{45}$.
Moreover
$\mathcal P_8/(T_{73}/R_8)=Q_8\times C_3^2$.
In the actual $[3]$ covering $E\setminus E[3]\to
E\setminus\{0\}$, all nine peripheral loops of $K_3$
are conjugates of $[a,b]$ or its inverse.  They map to the
central element $-1$ of $Q_8$.  Thus the degree-eight cover
has $36$ peripheral components and genus $19$; its Euler
characteristic is $-72$, in agreement with rank $73$.
This is the quaternion cover of the specified three-division
surface, with no replacement of its marked puncture set.

The complete producers and certificates are
\path{work/tower_code/modular_quaternion_surface.py},
\path{work/tower_code/modular_quaternion_class2.py},
and their matching \path{.json}, \path{.npz}, and
\path{.out} files.  The first producer additionally writes
an independent permutation representation and GAP check of
the linear base.  The class-two calculation itself uses
exact free-word substitution, finite-field row reduction,
and exhaustive finite Cayley enumeration only.
The independent root reproduction starts with copies of the
two programs and no precomputed JSON or NPZ input.  It gives
the same six surviving residues.  The separate verifier
\path{code/verify_quaternion_affine_census.py} then enumerates
all $3^8$ kernel vectors for every one of the $36$ candidates:
the six surviving fibres each have $81$ elements and the
other thirty are empty.  It also enumerates the commuting
matrices directly among all $12^4$ matrices, checks every
one of the $36864$ Schreier rows against the retained span,
and checks that the eight specified source rows generate it.
\end{proof}
